\documentclass[../main.tex]{subfiles}
\graphicspath{{\subfix{../img/}}}
\begin{document}
	\subsection{Introduzione}
	Il corso di Fondamenti di Elettronica si occuperà principalmente dell’utilizzo dei dispositivi elettronici realizzati con i semiconduttori.\\
	La principale differenza rispetto a quanto visto nei corsi di Elettrotecnica, consiste nel fatto che i dispositivi elettronici non hanno una relazione lineare tra corrente e tensione.\\
	I circuiti realizzabili con i dispositivi elettronici sono suddivisibili in due gruppi:
	\begin{itemize}
		\item Analogici
		\item Digitali
	\end{itemize}
	In questo corso tratteremo solo i circuiti digitali lasciando la trattazione dettagliata dei circuiti analogici a corsi più specifici.\\
	Uno schema generico del corso in oggetto è riportato in Figura \ref{fig:im-1}.
	\begin{center}
		\begin{figure}[h]
			\centering
			\includegraphics[scale=0.55]{1.png}
			\caption{Origini ed oggetto di studio del corso}
			\label{fig:im-1}
		\end{figure}
	\end{center}
	\subsection{Semiconduttori}
	\subsubsection{Premessa}
	Sappiamo che l’intensità di corrente ($J$) è legata al campo elettrico ($E$) dalla relazione:
	$$\vec{J} = \sigma \vec{E}$$
	Se la conducibilità ($\sigma$) è elevata, siamo in presenza di un \textbf{conduttore}; se $\sigma$ è ridotta siamo in presenza di un \textbf{isolante}.\\
	Possiamo dunque dire che i materiali \textbf{semiconduttori} sono materiali che si pongono a metà strada tra i conduttori (materiali ben disposti a condurre corrente sotto l’azione di un campo elettrico) e gli isolanti.\\\\
	La definizione appena data è però solo quantitativa. \\
	In realtà sappiamo che all’aumentare della temperatura, nei conduttori la conducibilità tende a diminuire mentre nei semiconduttori abbiamo il comportamento contrario: \textbf{all’aumento della temperatura, la conducibilità di un semiconduttore aumenta}.\\\\
	Dalla Fisica e dalla Chimica sappiamo inoltre che gli elettroni ruotano attorno al nucleo dell’atomo su un’orbita e con una velocità tali da controbilanciare esattamente la forza attrattiva del nucleo.\\
	Dall’orbita percorsa dall’elettrone possiamo allora associare ad ogni elettrone una energia (l’energia cinetica) ed una forza (la forza attrattiva del nucleo).\\
	In particolare, in un \textbf{sistema atomico}, sono permessi solo alcuni valori di energia (alcune orbite non sono cioè percorribili).\\
	Siccome gli elettroni hanno natura sia corpuscolare sia ondulatoria è anche necessario che al termine di un’orbita un elettrone abbia effettuato un multiplo intero di periodi ondulatori.
	\subsubsection{Conclusioni}
	Non tutti i valori energetici visualizzabili su un asse sono accessibili agli elettroni.\\
	In altre parole l’energia degli elettroni è quantizzata (principio di quantizzazione dell’energia).\\\\
	Ogni valore energetico consentito può essere assunto da uno ed un solo elettrone (principio di esclusione di \textbf{Pauli} \footnote{In realtà su ogni livello possono essere presenti fino a due elettroni caratterizzati da un momento di spin opposto.}).\\
	Dal principio di esclusione deduciamo che esistono livelli occupati e livelli liberi.\\
	In condizioni di quiete gli elettroni si posizionano spontaneamente sui livelli ad energia più bassa.\\
	È inoltre possibile determinare, tramite analisi statistiche, un livello energetico (detto livello energetico di \textbf{Fermi}, $E_F$) che discrimina tra i livelli completamente occupati ed i livelli liberi in condizione di quiete.
	\subsubsection{Sistemi Multiatomici}
	I livelli energetici rilevati nei sistemi monoatomici sono leggermente diversi dai livelli energetici rilevabili sui sistemi biatomici ma tutte le proprietà precedenti sono invariate. \\
	Avvicinando sempre più due atomi arriveremo alla fine a creare un legame tra i due detto legame covalente. \\
	Solitamente gli elettroni che formano il legame sono quelli con livelli energetici più elevati.\\\\
	Aumentando sempre più il numero degli atomi avremo fasce molto dense di livelli energetici distinti. \\
	Allora avremo una \textbf{struttura energetica a bande}. \\
	Una struttura energetica a bande è una struttura energetica in cui si alternano bande dense di livelli energetici permessi a bande proibite.\\\\
	Per muovere un elettrone allora sarà necessaria un’energia tale da spostarlo dal livello energetico in cui si trova ad un altro livello energetico permesso e libero. \\
	Allo zero assoluto l’elettrone al livello più prossimo al livello di Fermi sarà quello per il quale l’energia utilizzata sarà minima.\\\\
	Nei conduttori, il livello di Fermi è interno ad una banda permessa e, pertanto, per eccitare gli elettroni è necessaria poca energia.\\ 
	Negli isolanti, invece, il livello di Fermi è posto in una banda proibita e quindi è necessaria una grande quantità di energia per eccitare gli elettroni.\\\\
	Il semiconduttore ha dunque una struttura quantitativa uguale a quella di un isolante. \\
	La differenza è qualitativa: il gap tra due bande permesse è molto più piccolo rispetto allo stesso gap di un isolante.\\\\
	Sono semiconduttori quei materiali in cui, a temperatura ambiente, gli elettroni possono essere eccitati per il semplice effetto dello scambio termico con l’ambiente. \\
	Una volta eccitato l’elettrone si troverà in una zona con molti livelli energetici liberi in cui muoversi e condurre corrente.\\
	La Figura \ref{fig:im-2} mostra la struttura energetica di isolanti, semiconduttori e conduttori.
	\begin{center}
		\begin{figure}[h]
			\centering
			\includegraphics[scale=0.55]{2.png}
			\caption{Origini ed oggetto di studio del corso}
			\label{fig:im-2}
		\end{figure}
	\end{center}
	Un’altra differenza fondamentale è che a mano a mano che gli elettroni a livelli energetici più elevati si spostano ad energie superiori al livello di Fermi, i livelli da essi precedentemente occupati vengono via via liberati. \\
	Allora avremo due bande parzialmente occupate in cui sarà necessaria poca energia per eccitare altri elettroni.\\\\
	La tabella \hyperref[tab:table-1]{1} riassume quanto appena visto.
	\begin{table}[h]
		\centering
		\label{tab:table-1}
		\begin{tabular}{|l|c|c|c|}
			\hline
			& $E_F$ & Gap & Bande di conduzione \\ \hline
			\textbf{Conduttore} & Interno & Molto ridotto & Una \\ \hline
			\textbf{Semiconduttore} & Esterno & Ridotto & Due \\ \hline
			\textbf{Isolante} & Esterno & Elevato & Due \\ \hline
		\end{tabular}
		\caption{Caratteristiche principali dei conduttori, dei semiconduttori e degli isolanti}
	\end{table}
	La banda immediatamente superiore al livello di Fermi è detta \textbf{banda di conduzione} mentre la banda inferiore al livello di Fermi è detta è detta \textbf{banda di valenza}.\\\\
	I materiali semiconduttori sono tutti i materiali posti sulla quarta colonna della tavola periodica degli elementi. \\
	In particolare useremo il \textbf{Silicio (Si)} ed il \textbf{Germanio (Ge)} ed alcuni materiali composti quali l’arseniuro di gallio (GaAs) ed il fosfuro di indio (InP).
	\subsection{Il Silicio (Si)}
	Il silicio (simbolo Si; numero atomico 14) è un elemento della quarta colonna della tavola periodica degli elementi. \\
	La banda di valenza del silicio presenta quattro elettroni. \\
	Il silicio genera un reticolo cristallino cubico a facce centrate che è rappresentabile in due dimensioni come in Figura \ref{fig:im-3}.
	\begin{figure}[h]
		
		\begin{subfigure}{0.5\textwidth}
			\includegraphics[width=0.9\linewidth, height=6cm]{3.png} 
			\caption{Reticolo cristallino del silicio sviluppato su due dimensioni}
			\label{fig:im-3}
		\end{subfigure}
		\begin{subfigure}{0.5\textwidth}
			\includegraphics[width=0.9\linewidth, height=6cm]{4.png}
			\caption{Reticolo cristallino tridimensionale del silicio}
			\label{fig:im-4}
		\end{subfigure}
		
		\caption{}
		\label{fig:image1}
	\end{figure}
	In condizione di quiete la carica del cristallo di silicio è nulla in qualunque suo punto. \\
	Per tale motivo la densità locale di carica ($\rho$) è costante in qualsiasi punto del dispositivo.\\\\
	Quando il cristallo di silicio riceve energia sufficiente a consentire ad uno o più elettroni il salto dell’energy gap, il dispositivo resta globalmente neutro.\\\\
	Tuttavia a livello locale, in un punto del dispositivo avremo un elettrone in più mentre in un altro punto ne avremo uno in meno.\\\\
	A questo punto l’assenza di un elettrone in un punto del cristallo provoca una carica locale positiva. \\
	Questa carica positiva però attira a sé nuovi elettroni. \\
	Per compiere il salto da un livello occupato al livello liberato dal precedente elettrone serve un quantitativo di energia molto minore rispetto alla quantità necessaria per liberare il primo elettrone.\\\\
	Allora in un semiconduttore si hanno due meccanismi: una carica negativa si muove seguendo l’elettrone ed una carica positiva si muove seguendo l’assenza di un elettrone di legame. \\
	Poiché il trasporto di energia è dovuto a due portatori di carica di segno opposto, si parla di trasporto bipolare.\\\\
	Dal momento che la carica positiva è dovuta all’assenza di un elettrone si dice che essa è dovuta ad una lacuna (in inglese è detta hole, buco).
	\subsection{Elettroni e Lacune}
	Le lacune e gli elettroni vengono solitamente misurati in densità poiché è più importante conoscere quante lacune od elettroni vi sono in un volume. \\
	Indicheremo con $n$ il \textbf{numero di elettroni liberi }nella banda di conduzione per unità di volume (l’unità di misura sarà l’inverso di un volume); indicheremo invece con $p$ il \textbf{numero di lacune} nell’unità di volume.
	$$n = \frac{\text{elettroni liberi in conduzione}}{\text{volume } (cm^3)}$$
	$$p = \frac{\text{lacune}}{\text{volume } (cm^3)}$$
	Per quanto detto fino ad ora è evidente che gli elettroni e le lacune nascono contemporaneamente. \\
	Dal momento che esiste l’evento di generazione di elettrone e lacuna deve esistere, affinché il sistema rimanga stabile, anche un meccanismo contrario per il quale le coppie scompaiano. \\
	La scomparsa delle coppie avviene quando diminuendo la propria energia, l’elettrone trova una lacuna e vi si insedia.\\\\
	L’energia può essere ottenuta o ceduta in vari modi (tra i più comuni vi sono il calore, le radiazioni elettromagnetiche e la luce).\\\\
	Indichiamo con $G$ il tasso di generazione e con $R$ il tasso di ricombinazione.
	$$G = \frac{\text{coppie generate}}{\text{volume tempo } (cm^3 s)}$$
	$$R = \frac{\text{coppie ricombinate}}{\text{volume tempo } (cm^3 s)}$$
	In condizioni di stazionarietà deve essere $G = R$. Questa uguaglianza non indica che non vi siano coppie: significa semplicemente che per ogni coppia creata in un dato istante, una seconda coppia scompare.
	Possiamo inoltre vedere $p$ ed $n$ come funzione della temperatura:
	$$n = p = n_i(T) \implies p \cdot n = n_i^2$$
	Alla temperatura di riferimento di 300° Kelvin (circa 27° Celsius), abbiamo:
	$$n_i(300) \approx 10^{10} \text{ cm}^3$$
	Siccome in un centimetro cubo esistono circa $10^{23}$ atomi, la generazione di coppie interessa un atomo ogni $10^{13}$ atomi. Nonostante sia un valore così ridotto rispetto al numero totale di elettroni, esso è comunque sufficiente a far condurre il materiale.
\end{document}